\section{Geometric exclusions for a nonvanishing counterexample}
\label{sec:exclusions}

We now prepare the smooth nonvanishing step. Assume only the
lower-dimensional good-model hypothesis of
Assumption~\ref{mmp:lower-models}, and suppose that a smooth
projective \(n\)-fold \(X\), with \(n>0\), satisfies
\begin{equation}\label{ex:counterexample}
 K_X\ \text{is pseudo-effective},
 \qquad \kappa(X,K_X)=-\infty.
\end{equation}
We first restrict its moving proper subvarieties and positive canonical
currents, and then exclude covering curves of bounded genus and normalized
degree. Two further conclusions will handle the remaining jet
configurations: the reduced-boundary adjoint of a signed canonical
representative is big, and fixed finite covers admit no family of
birational maps whose graphs dominate their product. The argument takes
place over \(\C\), and the two displayed properties persist on smooth
projective birational models. Numerical dimension for a
pseudo-effective divisor means Nakayama's numerical dimension
\(\kappa_\sigma\); for a nef divisor it agrees with the intersection
definition used in Theorem~\ref{thm:signed}.

\subsection{The Albanese reduction and algebraic webs}

Subadditivity first shows that the Albanese map is constant. This
will let us pass from an effective numerical representative of a
rational divisor to an effective rational-linear representative.
We then construct the rational quotient generated by a moving
family of subvarieties.

\begin{lemma}\label{ex:irregularity}
Every smooth projective birational model of \(X\) has irregularity zero.
Consequently, on such a model, or on a projective \(\Q\)-factorial
terminal birational model, numerical equivalence of rational divisors
implies their \(\Q\)-linear equivalence.
\end{lemma}
\begin{proof}
If \(q(X)>0\), resolve the Stein factorization of the Albanese map to
obtain a morphism \(f:W\to B\) with connected fibers, with \(W,B\)
smooth projective and \(\dim B>0\). The map from \(B\) to a subvariety
of \(\Alb(X)\) is generically finite. Wedges of invariant one-forms
on the abelian variety therefore give a nonzero section of \(K_B\):
at the generic point, choose \(\dim B\) independent pulled-back
one-forms. Thus \(\kappa(B,K_B)\geq0\).

For a very general smooth fiber \(F\), the restriction of the
pseudo-effective class \(K_W\) is pseudo-effective. For example, restrict
a positive current representing it to almost every fiber, or use
effective approximations with arbitrarily small ample error.
Adjunction identifies this restriction with \(K_F\). If \(F\) has
positive dimension, the induction hypothesis gives
\(\kappa(F,K_F)\geq0\); for a point the same assertion holds by
convention. Applying Theorem~\ref{asm:iitaka} with both boundaries
empty gives \(\kappa(W,K_W)\geq0\), contradicting
\eqref{ex:counterexample}.

This is the only use of Theorem~\ref{asm:iitaka} in the proof.
Irregularity is a smooth birational invariant. Finally, when
\(\Pic^0=0\), a numerically trivial line bundle is torsion, since the
group of numerically trivial line bundles modulo \(\Pic^0\) is
finite. Clear denominators for rational divisors. On a terminal
\(\Q\)-factorial model, pull back to a smooth resolution and then
descend the rational linear equivalence.
\end{proof}

To handle subvarieties through very general points, parameterize
them in countably many Hilbert families, then resolve and stratify.
On a suitable open parameter space the domains form a smooth
projective family with integral fibers, and evaluation is dominant.
When each fiber maps generically finitely onto its image, general
ample cuts of the parameter space make total evaluation generically
finite while preserving dominance. We resolve and compactify this
total evaluation before using its ramification divisor.
Invariance of plurigenera for smooth projective families
\cite[Theorem~1]{Paun2007} permits the same construction for the
Iitaka fibers of general members: choose a divisible pluricanonical
system on the general member, spread it by base change, and resolve
its relative rational map. When the member has nonnegative Kodaira
dimension but is not of general type, the smooth general fibers
of this map have positive dimension. They form the new family
used below.

\begin{lemma}[Algebraic web quotient]\label{ex:web}
Let a smooth projective variety be dominated generically finitely by
the total space of a family with smooth integral projective general
fibers. On a dense regular open, form the distribution generated by
the tangent spaces of the fiber images, taking spans, saturation,
and Lie brackets. Its general leaves are dense opens of algebraic
subvarieties. Their closures are the general fibers of a rational
map, which has connected general fiber after Stein factorization.
\end{lemma}
\begin{proof}
Shrink to an open \(U\) where the evaluation has finitely many etale
sheets, the parameter map is smooth, and the generated involutive
distribution \(\mathcal F\) has constant rank. Its construction is
algebraic: spans and brackets stabilize at the generic point after
finitely many operations, and descend from the etale sheets. The
nonempty open parts of the integral parameter fibers are connected.
A chain step moves between two points in the image of one such
open fiber. We use only steps lying in \(U\), so each step remains
in one analytic leaf of \(\mathcal F\).

For \(x\in U\), endpoints of chains of at most \(k\) steps form a
constructible set, by the algebraic fiber-product construction and
Chevalley's theorem. Every endpoint lies in the analytic leaf
through \(x\). A locally closed smooth variety contained in an
analytic leaf has dimension at most \(\rk\mathcal F\). Indeed, in a
Frobenius chart the leaf meets at most countably many plaques: use
finite plaque chains in a countable foliation atlas. The transverse
coordinate map on any connected smooth piece then has countable
image and is constant.

There is therefore a maximal dimension among the closures of all
such endpoint loci; it is attained for some finite \(k\). The loci
are nested because we allow at most \(k\) steps, including the
identity step. Choose an irreducible component \(V\) of maximal
dimension and a dense open \(O\subset V\) of reachable points.
On each etale sheet, the parameter map defines the algebraic
relation consisting of pairs in the same fiber. Restrict its
first coordinate to the inverse image of \(O\), take the component
through the diagonal, and map the second coordinate to \(U\).
The closure of this image contains \(O\), hence \(V\), while its
points are reachable in at most \(k+1\) steps. Maximality forces
that closure to equal \(V\). At a general
diagonal point, smoothness of the parameter map identifies its
vertical tangent with the corresponding web tangent. Every web
tangent is consequently tangent to \(V\), and so is every bracket. Hence
\(\dim V\geq\rk\mathcal F\), while the reverse inequality was proved
above. A plaque is thus open in \(V\). The equations of \(V\),
pulled back to the connected immersed leaf, vanish on a nonempty
open and hence on the entire leaf. Its closure is exactly \(V\).

These invariant subvarieties of dimension \(\rk\mathcal F\) have countably many
Hilbert or Chow parameter spaces. Tangency on the regular locus is
an algebraic condition after stratification, so one such family
dominates. An invariant irreducible variety meeting \(U\) contains
the leaf through any of its points in \(U\): the tangent vector
fields preserve its reduced ideal, first on its smooth locus and
then everywhere by closure. A leaf closure of dimension \(\rk\mathcal F\) is
therefore unique through a general point. Assigning that closure
gives the desired rational map. Resolve its graph and take the
Stein factorization.
\end{proof}

For positive closed \((1,1)\)-currents we shall use this quotient in
the following way. If a current vanishes on almost every parameter
fiber away from a fixed removed divisor, it annihilates the
corresponding fiber tangents on a dense open. In submersion
coordinates, choose locally integrable plurisubharmonic potentials.
Fubini's theorem makes the pure fiber Hessian zero as a
distribution; positivity then makes the mixed coefficients in
fiber directions zero as well. A generically etale evaluation
transfers this assertion to the target. Annihilation passes to
brackets by the identities
\[
 \mathcal L_v T=d(\iota_vT)+\iota_v(dT),\qquad
 \iota_{[v,w]}T=\mathcal L_v(\iota_wT)-\iota_w(\mathcal L_vT).
\]
Pullbacks by dominant morphisms and restrictions to almost every
smooth parameter fiber are defined by the same local potentials.

\subsection{The general-type exclusion}

The next proposition turns the lower-dimensional good-model hypothesis
into a restriction on every moving proper subvariety. Its proof passes
from Iitaka fibers to the algebraic quotient of their tangent
directions and then back to nonvanishing.

\begin{proposition}\label{prop:general-type}
Every positive-dimensional proper subvariety through a very general
point of \(X\) is of general type on resolution. The assertion holds
also on every projective birational model.
\end{proposition}
\begin{proof}
Otherwise take a covering family of nongeneral-type subvarieties,
with smooth domains \(V\), and slice its parameters as above.
Restricting the ramification formula of its generically finite
total evaluation shows that \(K_V\) dominates the restriction of
the pulled-back \(K_X\) by an effective divisor. Thus \(K_V\) is
pseudo-effective. The induction hypothesis gives a good minimal
model of \(V\), so \(\kappa(V)\geq0\). Since \(V\) is not of general
type, its general Iitaka fibers have positive dimension and
Kodaira dimension zero. Spreading these fibers gives another
covering family with smooth domains \(G\), of dimension less than
\(n\), and \(\kappa(G)=0\). Each \(G\) has a good minimal model and
\(\kappa_\sigma(K_G)=0\).

To exploit this zero-Iitaka-dimension family, fix a positive current
\(T\) representing \(K_X\). After slicing and
resolving the total evaluation of the \(G\)'s, its restriction to
almost every parameter fiber, plus the effective restricted
ramification divisor, is a positive current in \(K_G\).
Every positive current in that class is the current of the fixed
canonical divisor of \(G\). To see this, let \(R\) be such a
current and take a common resolution
\(u:H\to G\), \(v:H\to G_{\min}\) with the good minimal model.
The current
\[
 u^*R+[K_H-u^*K_G]
\]
is positive and represents \(K_H\). Since \(K_{G_{\min}}\) is
torsion, this class is represented by an effective
\(v\)-exceptional divisor. Its intersection with the
\((\dim G-1)\)-st power of an ample pullback from \(G_{\min}\)
is zero. That form is strictly positive off the exceptional
locus, so the current is supported there. The support theorem
makes it divisorial, and independence of exceptional divisor
classes, by negativity, determines its coefficients. Pushing
forward by \(u\) determines \(R\) itself.
After shrinking the parameter space, these fixed canonical
divisors are the fixed divisor of a sufficiently divisible
relative pluricanonical system. Remove this divisor and the
excluded loci of the construction from the total family. Their
images under generically finite evaluation are proper algebraic
subsets of \(X\).

It follows that \(T\) annihilates the web distribution of the
\(G\)'s on a dense open. If that distribution has full rank, \(T\)
is supported on a proper algebraic set. The support theorem for
positive closed \((1,1)\)-currents expresses it as a finite
nonnegative real combination of prime divisors
\cite{Siu1974,Demailly2012}. Its rational cohomology class then has
a nonnegative rational representative: solve the finite rational
linear system for the coefficients on the same face of the
nonnegative orthant. Lemma~\ref{ex:irregularity} converts numerical
effectivity to \(\Q\)-linear effectivity, a contradiction.

Otherwise Lemma~\ref{ex:web} gives a rational quotient with general
smooth fiber \(F\) on a smooth resolved graph, where
\(0<\dim F<n\). Its canonical divisor is pseudo-effective:
restrict the pseudo-effective canonical class of the smooth graph
to a very general fiber, as in Lemma~\ref{ex:irregularity}.
The moving \(G\)'s still generate its tangent space at general
points. Indeed the quotient is constant on every general web
member, hence induces a rational map on their parameter space;
restricting to a general quotient fiber preserves dominance of
evaluation. A pluricanonical rational map of \(F\) is constant
along each such \(G\). Slice the parameters inside \(F\) to make
the total evaluation generically finite: ramification injects the
restricted pluricanonical sections into sections of a multiple of
\(K_G\), and these span a space of dimension at most one.
Fix a nonzero denominator section and work where it does not
vanish. Its restriction is nonzero on a general member of the
covering family, so all ratios are constant on that member. Their
differentials vanish on the web and its brackets, so the
pluricanonical map is constant on \(F\). Lower-dimensional
nonvanishing therefore gives \(\kappa(F)=0\), and its good minimal
model gives \(\kappa_\sigma(K_F)=0\).

The quotient has now reached the exact setting of the
numerical-zero-fiber reduction of
Gongyo--Lehmann \cite[Theorem~1.3 and Corollary~4.5]{GongyoLehmann2013}.
For a projective \(\Q\)-factorial klt rational pair with a
connected-fiber morphism and numerical dimension zero on the
general fiber, that theorem produces a klt pair on a smooth
birational base whose good-minimal-model existence is equivalent.
Here the total space is the smooth graph, the boundary is zero,
and the base has dimension less than \(n\). The required numerical
dimension is \(\kappa_\sigma\), which is zero by the good model of
\(F\); thus the numerical-dimension qualification, with the
corrected convention discussed in
\cite[Section~3 and Theorem~3.2]{Fujino2019Numerical},
causes no issue. The induction hypothesis supplies the base good
model, hence nonvanishing upstairs, contradicting
\eqref{ex:counterexample}.
\end{proof}

\begin{corollary}\label{ex:supporting}
Let \(Y\) be a projective \(\Q\)-factorial terminal birational model
of \(X\). If a rational divisor \(P\) has \(\kappa(Y,P)\geq1\),
then \(K_Y+cP\) is big for every rational \(c>0\).
Consequently, if \(\alpha\neq0\) is a supporting functional of the
pseudo-effective cone with \(\alpha(K_Y)=0\), then
\(\alpha(P)>0\).
\end{corollary}
\begin{proof}
Resolve a moving subsystem of a multiple of \(P\). If its map is
generically finite, \(P\) is big and the assertion follows from
pseudoeffectivity of \(K_Y\). Otherwise its general fiber is a
proper positive-dimensional subvariety through very general
points and is of general type by Proposition~\ref{prop:general-type}.
The canonical divisor of the smooth resolution is relatively big,
so adding a sufficiently large ample pullback from the image makes
it big. Since that pullback is bounded by a multiple of the
resolved moving subsystem, pushforward gives \(K_Y+aP\) big for
some \(a>0\). Convexity with the pseudo-effective \(K_Y\), and then
addition of the pseudo-effective \(P\), gives the assertion for
every \(c>0\). A nonzero nonnegative functional on a closed convex
cone is strictly positive on its interior. Applying it to
\(K_Y+cP\) proves the last assertion.
\end{proof}

\subsection{Valuations and positive currents}

The next two exclusions involve currents whose numerical intersection
gaps tend to zero. A limiting inequality alone is insufficient: we
need one family on which the gap is exactly zero. The ACC statement
below supplies that step for a fixed current and bounded discrepancies.

For a positive closed \((1,1)\)-current \(T\) on a smooth projective
variety \(V\), let \(\nu_E(T)\) denote the generic Lelong number of
its pullback at a divisorial valuation \(E\). The normalization is
such that a reduced smooth divisor has generic number one. We
write \(A_V(E)=a(E;V,0)\) for log discrepancy.

\begin{lemma}\label{ex:valuation-acc}
For a fixed \(T\) and a fixed \(M\), the set
\[
 \{\nu_E(T): A_V(E)\leq M\}
\]
satisfies the ascending chain condition.
\end{lemma}
\begin{proof}
We induct on the integer part of \(M\); discrepancies over a
smooth variety are positive integers, so the assertion is empty
when \(M<1\). Suppose that a strictly
increasing sequence exists, and discard initial terms so its
members exceed a fixed positive number \(c\).
Skoda integrability and change of variables give
\begin{equation}\label{ex:skoda-valuation}
 \nu_E(T)\leq A_V(E)\nu_x(T)
\end{equation}
at a general point \(x\) of the center. Indeed every exponent
smaller than \(1/\nu_x(T)\) is locally integrable, whereas
integrability after pullback imposes the corresponding
discrepancy bound. Hence all centers lie in the fixed proper Siu
locus \(\{\nu_x(T)\geq c/M\}\), which is algebraic by Siu
analyticity and projectivity \cite{Siu1974}. Passing to a
subsequence, all centers lie in one irreducible component of
this locus.

If, after passage to a subsequence, the centers lie in a fixed
codimension-at-least-two subvariety, principalize its ideal.
All lifted centers lie in the exceptional locus, where the
relative canonical divisor has positive integral order.
The discrepancy bound for the new smooth ambient space has
therefore decreased by at least one. Pull back \(T\) and apply
induction.

Otherwise the centers lie in a fixed prime divisor \(D\) of the
Siu locus. Write \(T=c_D[D]+T'\) with \(T'\) positive and with
zero generic Lelong number along \(D\). The integers
\(\ord_E(D)\) are bounded: the fixed effective divisor \(D\) has
positive log canonical threshold, and
\(\operatorname{lct}(D)\ord_E(D)\leq A_V(E)\).
Pass to a constant order. The residual numbers \(\nu_E(T')\)
are still strictly increasing, and after discarding a term have
a positive lower bound. Equation~\eqref{ex:skoda-valuation}
places their centers in a fixed Siu locus of \(T'\), which does
not contain \(D\). Their intersection with \(D\) has codimension
at least two, so the preceding case applies.
\end{proof}

To compare those Lelong numbers with movable algebraic systems, use
the following multiplier-ideal approximation
\cite[Theorems~5.11, 6.27, and 14.2]{Demailly2012}. On a smooth projective variety, if \(\{T\}\)
is a real algebraic class, one can choose integral line bundles
\(L_k\) such that \(c_1(L_k)-k\{T\}\) stays in a bounded,
uniformly sufficiently ample set and
\(L_k\otimes\mathcal J(kT)\) is globally generated. Round the
coefficients of \(k\{T\}\) in a fixed integral basis and add a
fixed sufficiently positive integral class. Nadel vanishing and
Castelnuovo--Mumford regularity give generation. For every
divisorial valuation,
\begin{equation}\label{ex:multiplier-order}
 \ord_E\mathcal J(kT)\leq k\nu_E(T).
\end{equation}
The local Bergman weight is bounded below by the original weight
up to a constant; this inequality persists after pullback and
gives \eqref{ex:multiplier-order}. These are statements for a
fixed current, not uniform assertions over all currents.

\begin{proposition}\label{prop:no-pullback}
Let \(Y\) be a fixed projective \(\Q\)-factorial terminal
birational model of \(X\), with \(K_Y\) nonbig. A positive current
in the pullback class of \(K_Y\) on a smooth resolution cannot,
on a dense regular Zariski open, annihilate the tangent spaces
of the general fibers of a rational fibration of positive
relative dimension.
\end{proposition}
\begin{proof}
Fix a current \(T\) as in the statement and suppose such a fibration
exists. Consider all dominant rational maps from \(Y\) with
connected general fiber of positive dimension for which the
pulled-back current annihilates the fiber tangents on a dense
regular open of a resolved graph. Choose one with smallest base
dimension \(b\). We will descend the current to this base and
find a covering family of curves on which the descended current
vanishes. Their quotient will give a member of the same collection
with smaller base dimension.
If \(b=0\), the current vanishes on
a dense open; the divisorial-current argument in
Proposition~\ref{prop:general-type} contradicts
\eqref{ex:counterexample}. Thus \(0<b<n\).
Resolve the graph and the space carrying the current:
\[
 \begin{tikzcd}
 W \arrow[r,"p"] \arrow[d,"h"'] & Y \\
 B &
 \end{tikzcd}
\]
Here \(W,B\) are smooth projective and \(h\) has connected
general fiber. All further modifications replace this graph and
base by birational models; they do not change the chosen rational
map or the function field of its base.
Flatten the main component over a modification of the base,
resolve that base, normalize the main transform, and then resolve
upstairs. Every vertical prime on the flat
transform maps to a divisor of the base: a prime over base
codimension \(c\) would have generic fiber dimension at least
\(n-b+c-1\), whereas all fibers have dimension at most \(n-b\).
A prime of \(W\) either maps onto a prime of the flat transform,
where this dimension count applies, or is exceptional over that
transform. In the latter case it is also exceptional over \(Y\),
because the transform maps birationally to \(Y\). Thus every
vertical divisor on \(W\) over base codimension at least two is
exceptional over \(Y\). This remains true after later resolutions
and base modifications: a prime dominating a divisor of \(Y\)
already has divisorial center on the old base, and its center on
a higher base model is the strict transform of that divisor.

For divisors and numerical classes put
\[
 T_B=p_*h^*.
\]
This strict pullback is well defined on numerical classes because
\(Y\) is \(\Q\)-factorial: numerical classes upstairs decompose
into pulled-back classes from \(Y\) and exceptional classes.
It preserves pseudoeffectivity. For a higher base model
\(r:B'\to B\), its analogue satisfies
\(T_{B'}r^*=T_B\), by computing on a common graph.
The operator \(T_{B'}\) kills divisors exceptional over \(B\):
otherwise a prime in such a pullback that dominates a divisor of \(Y\) would
have center of codimension at least two on \(B\), contrary
to the property just established.

\smallskip
\noindent\emph{Descent of the current.}
Let \(T\) also denote the pulled-back current on \(W\).
Shrink to a smooth open \(B^0\) so that \(h\) is smooth
proper, the removed set contains no vertical divisor
over \(B^0\), and the remaining open in every fiber is
nonempty and connected. Away from that fixed removed
algebraic set upstairs, \(T\) descends to a positive current
\(S^0\).
Indeed in submersion coordinates the horizontal coefficient
distributions are independent of fiber variables by closedness,
and the open parts of general fibers are connected. They
therefore glue to \(S^0\) on \(B^0\). The difference
\(T-h^*S^0\) is closed of order zero and supported on the
removed algebraic set. The support theorem expresses it as a
signed divisor sum; a closed order-zero \((1,1)\)-current
supported in codimension at least two is zero. Its horizontal coefficients are
nonnegative, since a pullback has zero generic Lelong number
along a horizontal divisor.

Subtract those finitely many global Siu components from
\(T\), obtaining a positive current \(T'\) that equals
\(h^*S^0\) on all of \(h^{-1}(B^0)\). Choose a K\"ahler form
\(\omega\) representing an ample algebraic class and put
\(d=n-b\). Fiber integration gives
\(h_*(\omega^d)=c>0\), constant on \(B^0\), by closedness.
The projection formula therefore gives
\[
 c^{-1}h_*(T'\wedge\omega^d)|_{B^0}=S^0.
\]
The left side defines a global positive closed extension.
Its class is real algebraic, since \(\{T'\}\) is the
algebraic class \(p^*K_Y\) minus real divisor classes and
algebraic intersection and pushforward preserve such
classes. Remove its divisorial parts along
\(B\setminus B^0\), and call the result \(S\).
On a dense open \(T=h^*S\). Their difference is again closed
of order zero, so the same support theorem leaves only a
signed divisor sum. Over \(B^0\), this difference consists
only of the horizontal components removed above. Its coefficient
at a prime dominating a base prime that meets \(B^0\) is therefore
zero. Along a base prime contained in \(B\setminus B^0\),
the current \(S\) has zero generic Lelong number by construction.
At a prime upstairs dominating it, the generic local map is a
transverse power map times a submersion. Local target balls of
a radius equal to a fixed power of the source radius give the
Lelong comparison, so \(h^*S\) also has zero generic number
there. Positivity of \(T\) then makes the coefficient of
\(T-h^*S\) nonnegative. Thus all coefficients at primes
dominating base primes are nonnegative.
The only possible negative coefficients lie over base
codimension at least two and are exceptional over \(Y\).
We have therefore obtained a positive closed current \(S\) in a
real algebraic class on \(B\), with \(T-h^*S\) a divisor current
whose only possible negative coefficients are \(p\)-exceptional.
Push this difference forward by \(p\). The exceptional terms
disappear and the remaining terms are effective, so
\begin{equation}\label{ex:descent-domination}
 K_Y-T_B\{S\}\quad\text{is pseudo-effective}.
\end{equation}
The horizontal components removed during descent contribute
effective divisors to the same difference.

\smallskip
\noindent\emph{A negative canonical direction on the base.}
Choose a nonzero supporting functional \(\alpha\) of the
pseudo-effective cone at \(K_Y\), and put
\(\eta=T_B^*\alpha\). The corresponding functionals \(\eta'\)
on higher smooth base models are movable classes by
pseudo-effective/movable duality \cite[Theorem~0.2]{BDPP2013}.
They annihilate base-exceptional divisors. Since \(S\) is positive,
\(\eta\cdot\{S\}\geq0\). Applying \(\alpha\) to
\eqref{ex:descent-domination} gives the reverse inequality,
because \(\alpha(K_Y)=0\). Hence \(\eta\cdot\{S\}=0\).
If \(c_D[D]\) is a positive Siu component of \(S\), both
\(S-c_D[D]\) and \([D]\) are positive. It follows that
\(\eta\cdot D=0\), or equivalently \(\alpha(T_BD)=0\).

Before applying multiplier approximation, we need to control the
divisorial part of this descended current. Only finitely many prime
divisors on \(B\) have positive
divisorial Lelong coefficient for \(S\). Otherwise their
nonzero strict pullbacks are effective divisors on \(Y\)
killed by \(\alpha\), with no common prime components.
Indeed, a nonexceptional prime of \(Y\) cannot map into the
intersection of two base primes.
Only finitely many base primes have zero strict pullback,
since their total pullbacks are supported in the finite
exceptional locus of \(p\). In the finite-dimensional rational
space of divisor classes, infinitely many remaining strict
pullbacks have a rational relation. This relation must have
both positive and negative coefficients: a nonzero effective
sum has positive intersection with an ample \((n-1)\)-fold
product on \(Y\), so cannot be numerically zero. The two sides
are thus nonzero effective rational divisors with no common
prime component. Lemma~\ref{ex:irregularity} makes them
\(\Q\)-linearly equivalent. Clearing denominators gives two
distinct linearly equivalent effective integral divisors and
hence a pencil, whose class is killed by \(\alpha\). This
contradicts Corollary~\ref{ex:supporting}.

We claim that
\begin{equation}\label{ex:negative-base}
 K_B\cdot\eta<0.
\end{equation}
The smooth general fiber \(F\) of \(h\) is of general type by
Proposition~\ref{prop:general-type}. The relative-positivity
theorem of Kov\'acs--Patakfalvi
\cite[Theorem~9.9]{KovacsPatakfalvi2017} applies to a log
canonical fiber space with smooth base, SNC total pair, and
log-general-type geometric generic fiber; for a rational base
divisor \(M\) with \(\kappa(M)\geq0\), it gives relative
subadditivity with the term \(h^*M\). Here both spaces are
smooth projective, the boundary is zero, and we take \(M=0\).
Invariance of plurigenera on the smooth locus identifies the
generic-fiber Kodaira dimension with that of \(F\). Thus
\begin{equation}\label{ex:relative-positive}
 \kappa(W,K_W-h^*K_B)\geq \kappa(F,K_F)=n-b>0.
\end{equation}
This is an established general-type-fiber theorem, not another
use or strengthening of Theorem~\ref{asm:iitaka}.
Choose compatible canonical divisors and set
\(P=K_Y-T_B(K_B)=p_*(K_W-h^*K_B)\).
For every sufficiently divisible \(m\), pushing forward an
effective divisor
\(\operatorname{div}_W(\varphi)+m(K_W-h^*K_B)\)
gives
\(\operatorname{div}_Y(\varphi)+mP\geq0\).
The resulting injection of sections preserves their ratios,
so \(\kappa(Y,P)\geq1\). Corollary~\ref{ex:supporting} now gives
\(0<\alpha(P)=-K_B\cdot\eta\), proving
\eqref{ex:negative-base}.

\smallskip
\noindent\emph{Rational curves and an exact zero gap.}
Apply the multiplier approximation to \(S\) on \(B\), and
principalize \(\mathcal J(kS)\) by \(r_k:B_k\to B\).
If \(F_k\) is its divisor, the class
\[
 P_k=\frac{r_k^*L_k-F_k}{k}
\]
is nef. Put \(\eta_k=T_{B_k}^*\alpha\). The class \(\eta_k\)
is movable and \(F_k\) is effective, so
\(0\leq P_k\cdot\eta_k\leq (L_k\cdot\eta)/k\).
Since \(\eta\cdot\{S\}=0\) and \(L_k-k\{S\}\) stays in a
bounded set, this gives the first estimate below. The second
uses the fact that \(\eta_k\) kills exceptional divisors:
\[
 0\leq P_k\cdot\eta_k=O(k^{-1}),\qquad
 K_{B_k}\cdot\eta_k=K_B\cdot\eta<0.
\]
Fix an ample divisor \(H\) on \(B\), and write
\(\delta=-K_B\cdot\eta>0\). On \(B_k\), choose an ample
divisor \(H_k\) and a sufficiently small rational
\(\epsilon_k>0\) so that the ample rational class
\[
 Q_k=P_k+k^{-1}r_k^*H+\epsilon_k H_k
\]
satisfies \(Q_k\cdot\eta_k=O(k^{-1})\). This is possible
because \(\eta_k\cdot r_k^*H=\eta\cdot H\) is fixed.
Approximate \(\eta_k\) by a positive sum
\(\gamma_k=\sum_{i=1}^{N_k} a_i[C_i]\), with \(a_i>0\),
of covering-curve classes,
closely enough that
\[
 Q_k\cdot\gamma_k\leq C/k,\qquad
 -K_{B_k}\cdot\gamma_k\geq\delta/2
\]
for a constant \(C\) independent of \(k\). Keep only terms
with \(K_{B_k}\cdot C_i<0\). Their total anticanonical
degree is at least \(\delta/2\), while their total
\(Q_k\)-degree is at most \(C/k\). Taking the weighted
average of the ratios therefore gives one covering curve
\(\Gamma_k\) with
\[
 \frac{Q_k\cdot\Gamma_k}{-K_{B_k}\cdot\Gamma_k}
 \leq\frac{2C}{\delta k}.
\]
The quantitative bend-and-break estimate
\cite[Theorem~5]{MiyaokaMori1986} gives a rational curve
through a general point of \(\Gamma_k\) with \(Q_k\)-degree
at most \(2\dim B\) times this ratio. Parameterization and
uncountability give a covering family of such rational curves,
denoted \(R_k\). Since \(Q_k-P_k-k^{-1}r_k^*H\) is ample
and both \(P_k\) and \(r_k^*H\) are nef, we obtain
\begin{equation}\label{ex:small-curves}
 P_k\cdot R_k=O(k^{-1}),\qquad
 H\cdot r_{k*}R_k=O(1).
\end{equation}
Shrink each parameter space so that the curve class and its
contacts with the exceptional divisors and the strict transforms
of the Siu divisors of \(S\) are constant. The general parametrized
member is free in characteristic zero. Consequently
\(K_{B_k}\cdot R_k\leq-2\). Put \(\bar R_k=r_{k*}R_k\).
Its bounded \(H\)-degree bounds \(K_B\cdot\bar R_k\), and
effectivity of the relative canonical divisor gives
\begin{equation}\label{ex:contact-bound}
 0\leq K_{B_k/B}\cdot R_k
   =K_{B_k}\cdot R_k-K_B\cdot r_{k*}R_k=O(1).
\end{equation}
In particular, the upper bound is independent of \(k\).
If an exceptional prime \(E\) has coefficient
\(a_E=A_B(E)-1\) in \(K_{B_k/B}\), a positive contact with
it contributes \(a_E(E\cdot R_k)\) to this bounded sum.
Both factors are positive integers. Thus the number of
exceptional primes met by \(R_k\), their log discrepancies
\(A_B(E)=a_E+1\), and their contact degrees are uniformly
bounded. Contacts with the finitely many strict transforms
of divisorial Lelong components of \(S\) are bounded by their
fixed degrees against \(\bar R_k\). Subtracting all these
divisorial parts from \(r_k^*S\) leaves a positive current.
Its restriction to almost every member of the family is
defined and positive, by the local-potential and Fubini
argument following Lemma~\ref{ex:web}. Its degree is the
first difference below; this proves that difference is
nonnegative. Using \eqref{ex:multiplier-order} and
\eqref{ex:small-curves} gives
\begin{equation}\label{ex:zero-gap}
 \begin{split}
 0&\leq \{S\}\cdot\bar R_k
          -\sum_E\nu_E(S)\,E\cdot R_k\\
  &\leq \{S\}\cdot\bar R_k-(F_k/k)\cdot R_k
   \leq O(k^{-1}).
 \end{split}
\end{equation}
The sum includes exceptional primes and those strict
divisorial components; terms with zero contact are irrelevant.
For the last bound, the expression involving \(F_k\) equals
\(P_k\cdot R_k+(\{S\}-L_k/k)\cdot\bar R_k\).
The first term is \(O(k^{-1})\) by \eqref{ex:small-curves};
the second has the same bound because \(L_k-k\{S\}\) stays
bounded and the \(H\)-degrees of \(\bar R_k\) are bounded.

Bounded ample degree gives only finitely many numerical
classes of the integral cycles \(\bar R_k\). Pass to one
class, so \(\{S\}\cdot\bar R_k\) is now one fixed real
number. The current \(S\) on \(B\) also remains fixed. By
Lemma~\ref{ex:valuation-acc}, the sums in
\eqref{ex:zero-gap} belong to an ACC set: there are a bounded
number of terms, their nonnegative integral multiplicities
are bounded, and all relevant discrepancies are bounded.
Finite sums of this kind preserve ACC, by successively
taking nonincreasing subsequences of their entries.
We can now replace the limiting gap by an exact equality: the
nonnegative gaps in \eqref{ex:zero-gap} must equal zero for some
covering families. Otherwise a sequence
tending to zero has a strictly decreasing positive
subsequence, giving a strictly increasing sequence of the
complementary sums.

For a family with zero gap, restrict the residual current to
almost every parameter fiber as above. Its integral on that
complete curve is zero. A positive current on a curve is a
positive measure, so zero integral forces it to vanish.
Off the removed divisors the residual current agrees with
\(r_k^*S\).
The discussion following Lemma~\ref{ex:web} and that lemma
itself yield a positive-relative-dimensional rational
quotient of \(B\) whose vertical directions are annihilated
by \(S\) on an open. Composing with \(h\) and taking the
Stein factorization gives a quotient of smaller base dimension
whose vertical directions are annihilated by \(T\). It belongs
to the collection minimized at the start, contradicting the
choice of \(b\).
\end{proof}

\subsection{Excluding normalized bounded curves}

We apply the same exact-zero mechanism to curves on the small
modifications of a fixed late model. The normalization by volume is
chosen so that its scale tends to zero, while the current and the
discrepancy lattice stay fixed.

Run a canonical LMMP on \(X\) with scaling of a fixed very ample
rational divisor \(A\). For rational \(t>0\), its positive-threshold
stages give terminal \(\Q\)-factorial models \(Y_t\) on which
\(K_t+tA_t\) is nef, big, and semiample
\cite{BCHM2010}. There are only finitely many divisorial
contractions, since each lowers Picard number. Fix a late model
\(Y\) after the last such contraction. All subsequent maps
\(Y\dashrightarrow Y_t\) are small, and are canonical
nonpositive birational maps. If the program terminates, use
the last model repeatedly. Define
\begin{equation}\label{ex:volume-scale}
 \mu_t=\vol(K_X+tA)^{1/n}.
\end{equation}
Since \(K_X\) is not big, \(\mu_t\to0\). In particular \(K_Y\)
is nonbig. The transform \(A_t\) is effective up to rational
linear equivalence; all intersections below use general
covering curves, which are not contained in a chosen such
representative.

\begin{lemma}[Exceptional current comparison]\label{ex:current-comparison}
Fix a positive current \(T\) on a smooth resolution in the
pullback class of \(K_Y\). On a common smooth resolution of
\(Y\dashrightarrow Y_t\), write
\[
 p^*K_Y=p_t^*K_t+J_t,\qquad J_t\geq0,
\]
where \(J_t\) is exceptional over \(Y_t\).
Then \(\nu_E(T)\geq\operatorname{coeff}_E J_t\) for every
prime on that resolution.
\end{lemma}
\begin{proof}
Pass to a common smooth model \(V\) also dominating the
fixed resolution carrying \(T\). Choose on the fixed
resolution a sufficiently positive line bundle \(H\).
For sufficiently divisible \(k\), multiplier approximation
gives a globally generated
\(\OO(kp^*K_Y+H)\otimes\mathcal J(kT)\).
Pull this system to \(V\). For the finite collection of primes
on the common resolution whose coefficients we are comparing,
choose a general member \(D_k\) of the underlying effective
divisor system. At every such prime \(E\), it has multiplicity
\(\ord_E\mathcal J(kT)\); global generation after removal
of the ideal ensures no additional generic vanishing.
Here the notation \(H\) also denotes its pullback to \(V\).

Put \(D'_k=p_{t*}D_k\) and \(H_t=p_{t*}H\). The first is
effective, and both are \(\Q\)-Cartier because \(Y_t\) is
\(\Q\)-factorial. Pushing the rational linear relation gives
\[
 D'_k\sim_\Q kK_t+H_t.
\]
Thus the exceptional divisor
\(D_k-p_t^*D'_k\) is rationally equivalent to
\[
 kJ_t+H-p_t^*H_t.
\]
The two exceptional divisors are equal: their difference is
numerically trivial and exceptional, so the negativity lemma
applied with both signs makes it zero. Since
\(p_t^*D'_k\) is effective,
\[
 \operatorname{coeff}_E D_k
 \geq k\operatorname{coeff}_E J_t
       +\operatorname{coeff}_E(H-p_t^*H_t).
 \]
The last coefficient is fixed for this model and \(t\).
Combining with \eqref{ex:multiplier-order}, dividing by
\(k\), and letting \(k\to\infty\) proves the assertion.
The same proof can be made on each common model, so the
comparison applies to all divisorial valuations needed
later. No uniformity in \(t\) of the fixed error is required.
\end{proof}

\begin{proposition}\label{prop:bounded-curves}
There are no sequences \(t_j\downarrow0\) and covering
families of curves on \(Y_{t_j}\), with smooth projective
domains \(C_j\), for which both
\[
 g(C_j),\qquad
 \frac{(K_{t_j}+t_jA_{t_j})\cdot C_j}{\mu_{t_j}}
\]
are bounded by a fixed constant. Degrees mean degrees on
the domains, or equivalently intersection with their
pushforward cycles.
\end{proposition}
\begin{proof}
Suppose such families exist. Fix once and for all a smooth
resolution \(p:W\to Y\) and a positive current \(T\) in
the class \(p^*K_Y\); pseudoeffectivity supplies this current.
Neither \(W\) nor \(T\) will depend on \(t\). Let \(g_0\)
be a common upper bound for the genera, and choose an integer
\(m>0\) such that \(mK_Y\) is Cartier.

We first compare the curves on these fixed models without
changing their domains. Slice the parameter space so that the
smooth total family has dimension \(n\), the parameter space
has dimension \(n-1\), and evaluation is generically finite.
It has rational maps to \(Y\), to \(W\), and to the common
models comparing \(Y\) and \(Y_t\). Each target is proper,
so each map extends at codimension-one points by the valuative
criterion. Its indeterminacy locus has dimension at most
\(n-2\), and its image cannot dominate the parameter space.
Remove these finitely many images from that space, then resolve
and compactify away from the remaining complete fibers.
Denote the resulting maps by
\(\rho_t:U_t\to W\) and \(q_t=p\circ\rho_t:U_t\to Y\).
The smooth general parameter fiber is still the original
\(C_t\), with the same genus. The current we use on this
varying space is always \(\rho_t^*T\).

For a prime \(Z\) of \(U_t\) exceptional over \(Y\), the
restricted valuation of its function field is
\(r_Z\ord_E\) for a divisorial valuation \(E\) over \(Y\).
The coefficient of \(Z\) in the ramification difference is
\begin{equation}\label{ex:ramification-cost}
 \operatorname{coeff}_Z(K_{U_t}-q_t^*K_Y)
       =r_Za(E;Y,0)-1.
\end{equation}
This follows by factoring through a model extracting \(E\)
and calculating the tame ramification of DVRs. It is
positive and bounded away from zero: \(Y\) is terminal,
so \(a(E;Y,0)>1\), and the index \(m\) puts these log
discrepancies in \(m^{-1}\Z\). In particular
\(a(E;Y,0)\geq1+1/m\). The
remaining ramification coefficients are nonnegative.
On the general parameter fiber,
\[
 K_{U_t}\cdot C_t=2g(C_t)-2,\qquad
 q_t^*K_Y\cdot C_t\geq0.
 \]
The latter inequality follows from pseudoeffectivity and
the covering property. Thus
\[
 0\leq (K_{U_t}-q_t^*K_Y)\cdot C_t\leq 2g_0-2.
\]
Each exceptional contact contributes
\((r_Za(E;Y,0)-1)(Z\cdot C_t)\) to this bounded sum.
The first factor is at least \(1/m\), and the second is
a positive integer. Formula~\eqref{ex:ramification-cost}
therefore bounds the number of contacted exceptional primes,
their log discrepancies, their ramification indices, and
their contact degrees. In particular the integers
\(r_Z(Z\cdot C_t)\) are uniformly bounded. Also
\[
 q_t^*K_Y\cdot C_t\in m^{-1}\Z\cap[0,2g_0-2],
\]
so these degrees take only finitely many values.

The ramification estimate has bounded the possible contacts. We next
compare their Lelong contributions with the canonical differences.
The generic Lelong number of \(\rho_t^*T\) at \(Z\) is
\(r_Z\nu_E(T)\). Extract \(E\), remove its generic
divisorial part, and use the transverse-power-map
comparison for the zero-generic-number remainder.
Lemma~\ref{ex:current-comparison} shows that these
coefficients dominate the coefficients of the pulled-back
canonical difference \(J_t\). Its support is exceptional
over \(Y\) as well as \(Y_t\), since the map between these
models is small. Subtracting all exceptional divisorial
parts of \(\rho_t^*T\) leaves a positive current. Shrink the
parameter space so that the relevant divisor contacts are
constant, and restrict this residual current to almost every
fiber. As in \eqref{ex:zero-gap}, its degree is the first
difference below and is nonnegative. Consequently
\begin{equation}\label{ex:bounded-gap}
 \begin{split}
 0&\leq q_t^*K_Y\cdot C_t
       -\sum_{Z\ {\rm exceptional}}
                r_Z\nu_E(T)\,Z\cdot C_t\\
  &\leq K_t\cdot C_t
   \leq (K_t+tA_t)\cdot C_t
   =O(\mu_t).
 \end{split}
\end{equation}
All divisor contacts used here are nonnegative, since the
curves cover the total space.

On the fixed smooth resolution \(p:W\to Y\), one has
\[
 a(E;W,0)=a(E;Y,0)
       -\ord_E(K_W-p^*K_Y)\leq a(E;Y,0),
\]
because the relative canonical divisor is effective.
Thus Lemma~\ref{ex:valuation-acc} applies to the one fixed
current \(T\) on \(W\), with a fixed discrepancy bound.
The sums in \eqref{ex:bounded-gap} form an ACC set: their
Lelong numbers belong to this ACC set, and both the number
of terms and their integral multiplicities
\(r_Z(Z\cdot C_t)\) are bounded. Pass to a fixed value of
\(q_t^*K_Y\cdot C_t\). Since \(\mu_t\to0\), the
exact-zero argument of \eqref{ex:zero-gap} gives a family
for which the left gap is zero.

For that family, the restriction of the residual positive
current to almost every parameter fiber is a positive
measure of total mass zero, so it vanishes. On the dense
open away from the removed exceptional divisors it agrees
with \(\rho_t^*T\). Apply Lemma~\ref{ex:web} to the family
evaluated by \(\rho_t:U_t\to W\). Its curve tangents generate
a rational fibration of positive relative dimension on the
smooth projective variety \(W\), and \(T\) annihilates the
fiber tangents on a dense open. Transport this rational
fibration through \(p:W\to Y\) on their common isomorphism
open.
This contradicts Proposition~\ref{prop:no-pullback}.
\end{proof}

\subsection{Special termination for signed representatives}

To turn a signed canonical representative into a big logarithmic
adjoint, we will run a program whose negative rays meet the reduced
boundary. The following special-termination statement uses only the
lower-dimensional hypothesis; its proof keeps the programs on the
boundary separate from the ambient one.

\begin{proposition}[Special termination over a point]
\label{prop:special-termination}
Assume Assumption~\ref{mmp:lower-models}.
Let \((V,\Delta)\) be a projective \(\Q\)-factorial dlt
\(n\)-fold with rational boundary and pseudo-effective adjoint.
Choose an effective ample rational divisor \(A\) such that
\((V,\Delta+A)\) is dlt and \(K_V+\Delta+A\) is nef.
The LMMP for \(K_V+\Delta\) with scaling of \(A\) has special
termination: if the program is infinite, its flipping loci
are eventually disjoint from the round-down of the boundary.

In particular, if at every stage the adjoint is rationally linearly
equivalent to a signed divisor supported on that round-down, the
program terminates at a log minimal model.
\end{proposition}

\begin{proof}
In an infinite ample-scaling program the scaling limit is zero:
a positive limit is covered by the termination theorem for a
dlt pair with an ample summand in the scaling divisor
\cite[Theorem~1.9(ii), equivalently Theorem~4.1(ii)]
{Birkar2012Special}. Discard the finitely many divisorial
contractions and write the flips as \(V_i\dashrightarrow V_{i+1}/Z_i\),
with scaling numbers \(\lambda_i>0\) tending to zero.
Fix a component \(S_1\) of \(\lfloor\Delta_1\rfloor\), write
\(S_i\) for its normal strict transform, and let \(T_i\) be the
normalization of its image in \(Z_i\). Since the ambient contraction
is small, \(S_i\to T_i\) is projective birational. The standard
discrepancy argument makes \(S_i\dashrightarrow S_{i+1}\)
an isomorphism in codimension one after finitely many steps
\cite[proof of Lemma~3.6]{Birkar2010}.

To identify the lower-dimensional input, consider the relative
program induced on this one component of the floor.
Take a small projective \(\Q\)-factorialization
\(p_i:S'_i\to S_i\), and put
\[
 D_i=K_{S'_i}+B_i=p_i^*((K_{V_i}+\Delta_i)|_{S_i}),
 \qquad C_i=p_i^*(A_i|_{S_i}).
\]
Adjunction gives a dlt pair \((S'_i,B_i)\), with \(B_i,C_i\geq0\),
and \((S'_i,B_i+\lambda_i C_i)\) is lc. The scaling divisor
has no floor component, so its restriction is effective.
The number \(\lambda_i\) is rational, being the ratio of
intersections of rational divisors on the contracted rational curve.
The globally nef divisor
\(L_i=K_{V_i}+\Delta_i+\lambda_i A_i\) descends rationally
linearly across \(V_i\to Z_i\). Indeed, for sufficiently small
rational \(\delta>0\), the pair
\((V_i,(1-\delta)\Delta_i)\) is klt and its negative adjoint
is ample over \(Z_i\). Relative basepoint freeness applies to a
Cartier multiple of \(L_i\); its numerical triviality and the
connected fibers give descent. Restricting and pulling back yields
\[
 D_i+\lambda_i C_i\sim_{\Q}0/T_i.
\]
The descended divisor on $T_i$ is nef: its pullback to $S'_i$
is the pullback of the restriction of the globally nef $L_i$.

By \cite[Theorem~1.1(3)]{Birkar2012Special}, a \(D_i\)-LMMP
over \(T_i\) with scaling of a fresh ample divisor terminates.
The theorem applies to the effective rational lc boundary
\(B_i+\lambda_i C_i\), with \(\lambda_i C_i\) rational Cartier
and the driving pair \(\Q\)-factorial dlt. Since the base map
is birational, the endpoint is a log minimal model, not a Mori
fiber space. Comparison with the relatively ample model
\(S_{i+1}\) identifies this endpoint with a small
\(\Q\)-factorialization \(S'_{i+1}\), as in
\cite[Remarks~2.9--2.10]{Birkar2012Special}.

The next point is that these finite relative programs concatenate
into an absolute program with scaling of the transforms of \(C_1\).
Throughout the \(i\)-th piece, \(D_i+\lambda_i C_i\) remains
the pullback of a nef divisor on \(T_i\), hence is globally nef.
Here $D_i$ and $C_i$ also denote their transforms along this
relative program. A contracted negative ray $R$ satisfies
\[
 (D_i+\lambda_i C_i)\cdot R=0,
 \qquad D_i\cdot R<0.
\]
It follows that $C_i\cdot R>0$. Every smaller coefficient of
$C_i$ leaves this ray negative, while $D_i+\lambda_i C_i$ is
globally nef. Thus the absolute scaling threshold at that step
is exactly $\lambda_i$.
The relative cone is a face of the absolute cone, cut out by
the pullback of an ample divisor on the projective base \(T_i\);
its extremal rays are therefore absolute extremal rays.
Thus the pieces concatenate as asserted in the cited remarks.
Rescaling the initial scaling divisor by \(\lambda_1\), if
necessary, makes its sum with the initial boundary lc and nef.

If the concatenation were infinite, its positive scaling numbers
would tend to zero. Fix a late scaling value $\lambda_j$.
On every earlier ray $R$ contracted at threshold $\lambda_i$,
the calculation above gives
\[
 (D_i+\lambda_j C_i)\cdot R
   =(\lambda_j-\lambda_i)C_i\cdot R\le0.
\]
Thus all preceding steps are nonpositive for the adjoint with
this fixed coefficient $\lambda_j$. On a common resolution of
the initial and the $j$-th models, the pullback of
$D_1+\lambda_j C_1$ equals the pullback of the late nef divisor
$D_j+\lambda_j C_j$ plus an effective comparison divisor.
Push this equality down to the initial model. The pushforward
of the nef pullback is pseudo-effective, as is the effective
comparison term, so $D_1+\lambda_j C_1$ is pseudo-effective.
Closedness of the pseudo-effective cone and $\lambda_j\to0$
show that $D_1$ is pseudo-effective.
Assumption~\ref{mmp:lower-models} therefore gives an absolute
log minimal model for that initial pair. The concatenation
terminates by \cite[Theorem~1.9(iii)]{Birkar2012Special}, since
its zero limit is never attained. The floor-component argument
of \cite[Lemma~3.6]{Birkar2010} now gives special termination.
The only good-model input in this argument concerns the absolute
pair on the lower-dimensional floor component. In particular it
does not require an effective representative in dimension \(n\).

For the final assertion, a curve disjoint from the round-down
has degree zero against any signed divisor supported there.
It therefore cannot generate an adjoint-negative flipping ray.
Special termination rules out all sufficiently late flips.
There are only finitely many divisorial contractions, since
each lowers the Picard number, and pseudo-effectivity excludes
a Mori fiber space as the endpoint. Thus the endpoint is nef.
\end{proof}


\subsection{The signed alternative}

The jet construction in Section~\ref{sec:jets} will produce signed
rational representatives of \(K_t\). The following result turns
such a representative into a big logarithmic adjoint on a resolution.

\begin{proposition}\label{prop:signed-alternative}
On any model \(Y_t\), let a signed rational divisor
represent \(K_t\) up to \(\Q\)-linear equivalence. On a
log resolution \(p:W\to Y_t\), let \(D\) be the reduced
SNC divisor consisting of its strict support and all
exceptional divisors. Then \(K_W+D\) is big.
\end{proposition}
\begin{proof}
Put \(P=K_W+D\), and suppose it is not big. It cannot
have Iitaka dimension at least one: by
Corollary~\ref{ex:supporting}, \(K_W+cP\) would be big,
and
\((1+c)P=(K_W+cP)+D\) would then be big as well.
Hence \(\kappa(W,P)\leq0\).

Run the dlt LMMP for \(P\) with ample scaling.
The divisor \(P\) has a signed rational representative
supported on the floor \(D\), and its transforms retain
that property. Every negative flip must meet the floor:
on its complement the representing rational section is
nowhere vanishing, so the adjoint has zero degree on
every complete curve there. By
Proposition~\ref{prop:special-termination}, an infinite
sequence would eventually have all flips disjoint from
the floor, a contradiction. Divisorial contractions
are finite in number. Thus the program terminates at
a \(\Q\)-factorial dlt log minimal model
\((Z,D_Z)\).

The boundary is reduced, \(K_Z\) is pseudo-effective
as the birational pushforward of \(K_W\), and the nef
adjoint has a signed rational representative supported
on \(D_Z\). Its restriction to \(D_Z\) is semiample
by Proposition~\ref{prop:boundary-restriction}.
Apply Theorem~\ref{thm:signed} with
\(L=K_Z+D_Z\) and \(c=1\), so that the required
pseudo-effective class \(L-cD_Z\) is \(K_Z\).
It gives \(\kappa(Z,L)=\nu(Z,L)\geq0\).
The Iitaka dimension is at most zero, as proved above,
so both dimensions are zero. Intersecting
\(K_Z+D_Z\equiv0\) with an ample
\((n-1)\)-fold product, and using pseudoeffectivity of
\(K_Z\), forces \(D_Z=0\). The signed relation then
gives \(K_Z\sim_\Q0\).

The signed representative of \(K_W\) is supported on
\(D\): pull back the given representative of \(K_t\)
and add the relative canonical divisor. Since \(D_Z=0\),
every component of this support is contracted over \(Z\).
On a common resolution \(a:V\to W\), \(b:V\to Z\),
we therefore have \(a^*K_W\sim_\Q E\) for a signed
\(b\)-exceptional divisor \(E\). The class \(a^*K_W\)
is pseudo-effective. A positive current in this class has
zero intersection with the \((n-1)\)-st power of an ample
pullback from \(Z\), so is supported on the exceptional
locus. The support theorem makes it effective divisorial.
Independence of exceptional divisor classes, by negativity,
identifies its coefficients with those of \(E\); hence
\(E\geq0\). Pushing forward by \(a\) shows that
\(K_W\) is \(\Q\)-linearly effective,
contradicting \eqref{ex:counterexample}.
\end{proof}

\subsection{Birational families between fixed covers}

The last obstruction concerns fixed finite covers of a birational model
of the counterexample. If birational maps between two such covers form
an algebraic family whose graphs dominate their product, then one fixed
cover is birational to an abelian variety. The positive canonical current
rules this out.

\begin{proposition}[Fixed-cover obstruction]\label{ex:fixed-cover}
Let \(Y\) be a projective \(\Q\)-factorial terminal birational model
of the counterexample \(X\) in Equation~\eqref{ex:counterexample}, and let
\(T_i\to Y\), \(i=1,2\), be fixed finite surjective morphisms from normal
integral projective varieties.
There is no integral parameter variety \(A_0\) and closed subvariety
\[
 \Gamma\subset A_0\times T_1\times T_2
\]
flat over \(A_0\), with each fiber the integral graph of a birational
map \(T_1\dashrightarrow T_2\), for which the projection
\(\Gamma\to T_1\times T_2\) is dominant.
\end{proposition}

\begin{proof}
Suppose such a family exists, and denote its birational maps by
\(\phi_a:T_1\dashrightarrow T_2\).
Choose one map \(\phi_{a_0}\) in this family and use it to identify
\(T_2\) birationally with \(T_1\). The maps
\(g_a=\phi_{a_0}^{-1}\circ\phi_a\) are then birational selfmaps of
\(T_1\). The evaluation \((a,u)\mapsto(u,g_a(u))\) is dominant onto
\(T_1\times T_1\); in particular, the images of a general point under
these selfmaps are dense.

We use the birational automorphism group to identify this fixed cover.
The cover is non-uniruled, being generically finite over the
non-uniruled \(Y\).  Hanamura's non-uniruled birational-group
theorem gives a smooth projective birational model for which the
reduced birational group is a group scheme, locally of finite type,
and its identity component is an abelian variety of regular automorphisms
\cite[Theorems~2.1--2.2]{Hanamura1988};
see also \cite[Proposition~3.7 and Theorems~3.8--3.9]{Blanc2017}.
After conjugating the maps to that model, shrink the irreducible
parameter variety anew so that their graph closures are flat
with integral fibers and both graph projections remain birational on
every fiber. This is a family in the flat graph functor of
\cite[Definition~3.4]{Blanc2017}, which is represented by the
birational scheme \cite[Proposition~3.7]{Blanc2017}. It therefore
defines a morphism to that scheme;
because the parameter variety is reduced, this morphism factors
through its reduction. Its connected image lies in a single component,
hence in a translate of the identity component. Translate the family
by one member so that it lies in that identity component; this
preserves dominance of evaluation. Thus the identity component, an
abelian variety, has a dense orbit. That orbit is the quotient by the
stabilizer of a general point. A quotient of an abelian variety is
again an abelian variety, so the fixed cover is birational to an
abelian variety.

It remains to see why this abelian cover contradicts the original
counterexample.  Resolve the generically finite rational map from
that abelian variety to \(X\):
on a smooth model \(U\), ramification gives
\[
 K_U=f^*K_X+R,\qquad R\ge0.
\]
The canonical class of \(U\) also has the effective representative
\(E_A\) exceptional over the abelian variety: choosing a nowhere-zero
top form on that variety gives \(K_U\sim E_A\). This second choice
of representative need not be the canonical divisor compatible with
the displayed ramification formula. Let \(T\) be a positive current in \(K_X\).
The current \(f^*T+[R]\) is positive and represents \(K_U\).
Let \(\alpha\) be an ample class on the abelian variety. The
intersection of \(K_U\) with the \((n-1)\)-st power of its pullback
is zero, since \(K_U\sim E_A\) and \(E_A\) is exceptional.
The positive current \(f^*T+[R]\) therefore has zero mass against
that power. On the open where the birational map to the abelian
variety is an isomorphism, the pulled-back ample form is strictly
positive, so the current vanishes there. Its support is consequently
exceptional over the abelian variety. The positive summand \(f^*T\)
is supported there as well, and the support theorem makes it divisorial.
Push forward this summand alone: \(f_*(f^*T)\) is a positive
divisorial current whose class is \(\deg(f)c_1(K_X)\) by the
projection formula.  Rationality of the numerical class supplies
a rational effective representative, and irregularity zero turns
numerical effectivity into nonvanishing, as in
Lemma~\ref{ex:irregularity}.  This contradicts
\(\kappa(X,K_X)=-\infty\).
\end{proof}
